\templateentry
{Hashing}
{\texttt{H()} é neutro; \texttt{H(x)} tem um símbolo. \texttt{a.concat(b)} retorna $ab$. Módulo $2^{61}-1$, base aleatória compartilhada; $O(1)$ por concatenação.}
{\texttt{H}: \texttt{H()} é vazio/neutro; \texttt{H(x)} tem um símbolo. \texttt{a.concat(b)} retorna $ab$ sem alterar os operandos; igualdade compara hash e comprimento. Base aleatória compartilhada, módulo $2^{61}-1$; colisões possíveis. Concatenação $O(1)$; cada hash guarda $\mathtt{pw}=\mathtt{base}^{\mathtt{cnt}}$.}
{strings/hashing.cpp}

\templateentry
{KMP}
{\texttt{kmp(s)} retorna a prefix function em $O(n)$.}
{\texttt{kmp(s)} retorna a prefix function em $O(n)$.}
{strings/kmp.cpp}

\templateentry
{Z-function}
{\texttt{zFunc(s)}: \texttt{z[i]} = LCP de \texttt{s} com \texttt{s[i..n)}; \texttt{z[0]=n}. Aceita string vazia. $O(n)$.}
{\texttt{zFunc(s)}: \texttt{z[i]} = LCP de \texttt{s} com \texttt{s[i..n)}; \texttt{z[0]=n}. Aceita string vazia. $O(n)$.}
{strings/z-function.cpp}

\templateentry
{Aho-Corasick}
{\texttt{AC}: adicione padrões \texttt{a--z} com \texttt{add}; chame \texttt{build} uma vez, depois \texttt{next(v,c)}. \texttt{add} retorna nó terminal; ocorrências exigem propagação pelos links. Construção $O(26|\mathtt{AC}|)$; percurso $O(|\mathtt{texto}|)$.}
{\texttt{AC}: adicione padrões \texttt{a--z} com \texttt{add}; chame \texttt{build} uma vez, depois \texttt{next(v,c)}. \texttt{add} retorna nó terminal; ocorrências exigem propagação pelos links. Construção $O(26|\mathtt{AC}|)$; percurso $O(|\mathtt{texto}|)$.}
{strings/aho-corasick.cpp}

\templateentry
{Array de sufixos e LCP}
{\texttt{sufArr(s)} retorna \texttt{\{sa,lcp\}}: sem sufixo vazio; \texttt{lcp[i]} entre \texttt{sa[i-1]} e \texttt{sa[i]}, \texttt{lcp[0]=0}. Requer $n>0$; $O(n\log n)$.}
{\texttt{sufArr(s)} retorna \texttt{\{sa,lcp\}}: sem sufixo vazio; \texttt{lcp[i]} entre \texttt{sa[i-1]} e \texttt{sa[i]}, \texttt{lcp[0]=0}. Requer $n>0$; $O(n\log n)$.}
{strings/suffix-array.cpp}

\templateentry
{Árvore de sufixos}
{\texttt{ST(s,sa,lcp)}: $O(n)$.}
{\texttt{ST(s,sa,lcp)}: árvore de sufixos em $O(n)$. Nós em \texttt{t}: \texttt{[l,r)} rotula a aresta do pai \texttt{p}; \texttt{dep} guarda profundidade em caracteres. \texttt{suf>=0} marca o fim do sufixo iniciado em \texttt{suf}; sem terminador único, pode ocorrer em nó interno.}
{strings/suffix-tree.cpp}

\templateentry
{Autômato de sufixos}
{\texttt{SAM sam}; \texttt{extend(c)} para letras \texttt{a--z}. \texttt{ord()} ordena estados por \texttt{len} crescente; propague \texttt{cnt} pelos links em ordem inversa. \texttt{t[v].end}: fim 0-based de uma ocorrência; comprimento $L$ dá intervalo $[\mathtt{end}-L+1,\mathtt{end}+1)$. Construção $O(n)$.}
{\texttt{SAM sam}; \texttt{extend(c)} para letras \texttt{a--z}. \texttt{ord()} ordena estados por \texttt{len} crescente; propague \texttt{cnt} pelos links em ordem inversa. \texttt{t[v].end}: fim 0-based de uma ocorrência; comprimento $L$ dá intervalo $[\mathtt{end}-L+1,\mathtt{end}+1)$. Construção $O(n)$.}
{strings/suffix-automaton.cpp}

\ifextended

\templateentry
{Maior substring comum}
{}
{\texttt{lcs(a,b)}: letras \texttt{a--z}; $O(|a|+|b|)$.}
{strings/longest-common-substring.cpp}

\fi

\templateentry
{Manacher}
{\texttt{manacher(s)}: \texttt{\{d1,d2\}}. Comprimentos $2\mathtt{d1}[i]-1$ (centro $i$) e $2\mathtt{d2}[i]$ (centro entre $i-1,i$). $O(n)$.}
{\texttt{manacher(s)}: \texttt{\{d1,d2\}}. Comprimentos $2\mathtt{d1}[i]-1$ (centro $i$) e $2\mathtt{d2}[i]$ (centro entre $i-1,i$). $O(n)$.}
{strings/manacher.cpp}

\ifextended
\templateentry
{Árvore de palíndromos (eertree)}
{}
{\texttt{add(c)} acrescenta uma letra \texttt{a--z} e retorna o nó do maior sufixo palíndromo atual. Nós $v\ge2$ representam palíndromos distintos; 0 e 1 são raízes de comprimentos $-1$ e 0. \texttt{link}: maior sufixo palíndromo próprio; \texttt{end}: fim de uma ocorrência, 0-based. \texttt{occurrences()[v]} conta todas as ocorrências, inclusive sobrepostas; contagem não destrutiva. \texttt{cnt} conta apenas vezes como maior sufixo. Construção $O(n)$ para alfabeto fixo; contagem $O(n)$, memória $O(26n)$.}
{strings/palindromic-tree.cpp}

% https://arxiv.org/abs/1506.04862
\textbf{Séries de sufixos palíndromos.}
Para estender a eertree, mantenha $d[v]=\mathtt{len}[v]-\mathtt{len}[\mathtt{link}[v]]$ nos nós não vazios. Agrupe nós consecutivos de mesmo $d$ na cadeia de suffix links em séries aritméticas: há $O(\log n)$ por prefixo. O \emph{series link} pula o grupo; com $u=\mathtt{link}[v]$,
\[
\mathtt{sl}[v]=\begin{cases}\mathtt{sl}[u],&d[v]=d[u],\\u,&d[v]\ne d[u].\end{cases}
\]
Inicialize $d=0$ e $\mathtt{sl}=1$ nas duas raízes do snippet; calcule os campos ao criar cada nó. Parta de \texttt{last} e siga $\mathtt{sl}$ enquanto $\mathtt{len}[v]>0$: a série de $v$ tem comprimentos de $\mathtt{len}[v]$ até $\mathtt{len}[\mathtt{sl}[v]]+d[v]$, em passos de $-d[v]$. Para acelerar DP de partição em palíndromos, agregue as transições de cada série; percorrer todos seus nós ainda pode custar $O(n)$ por prefixo.
\fi

\templateentry
{Decomposição de Lyndon}
{\texttt{lyndon(s)}: fronteiras $p=(0,\ldots,n)$ da fatoração; fatores $s[p_i,p_{i+1})$ em ordem lexicográfica não crescente. String vazia retorna \texttt{\{0\}}. $O(n)$.}
{\texttt{lyndon(s)}: fronteiras $p=(0,\ldots,n)$ da fatoração; fatores $s[p_i,p_{i+1})$ em ordem lexicográfica não crescente. String vazia retorna \texttt{\{0\}}. $O(n)$.}
{strings/lyndon.cpp}

\templateentry
{Encontrar runs em uma string}
{\texttt{getRuns(s)}: triplas \texttt{\{p,l,r\}} de repetições maximais em $[l,r)$, período mínimo $p$, $r-l\ge2p$. $O(n\log^2 n)$ nesta implementação.}
{\texttt{getRuns(s)}: triplas \texttt{\{p,l,r\}} de repetições maximais em $[l,r)$, período mínimo $p$, $r-l\ge2p$. $O(n\log^2 n)$ nesta implementação.}
{strings/runs.cpp}
